\scanpage{28}
% loc 1–3, continuation at upper edge
$\unsure{\widetilde A|\check A}$-module structure on $M$ such that for every $x\in M$, $\xi\mathrel{\tikz[baseline=-.3ex,x=1em,y=1em]{\draw[line width=.45pt] (0,-.13)--(0,.13) (0,0)--(1,0)--(1,.14);}}\xi x$ is continuous (i.e. vanishes on one of the $J_{\unsure{i|1|\prime}}\subset\check A\ \ldotp\ldotp\ldotp\ldotp\ \unsure{]|\exists}$. 
% loc 3–6, Likewise, we see
Likewise, we see that giving \deleted{\unsure{s}}a coalgebra structure \tikz[baseline=(s.base)]{\node[inner sep=1pt](s){on the};\node[anchor=south,inner sep=1pt] at ([yshift=6pt]s.north){\inserted{(\unsure{associative unital})}};\draw ([xshift=-4pt,yshift=-2pt]s.south west) to[out=75,in=180] ([yshift=5pt]s.north west) to[out=10,in=200] ([xshift=5pt,yshift=10pt]s.north east);} free $k$-module $A$ amounts to giving a \deleted{operation} $k$-algebra structure (associative unital) on $\check A$ which is \uline{continuous}
% loc 7–8, for if; struck line and formula
[for if $\overset{\inserted{\unsure{k^{\Sigma}|k^{I}}}}{\overset{\tikz[baseline=-.4ex]{\node[rotate=90,inner sep=0pt]{\(\scriptscriptstyle\simeq\)};}}{P}},\overset{\inserted{k^{J}}}{\overset{\tikz[baseline=-.4ex]{\node[rotate=90,inner sep=0pt]{\(\scriptscriptstyle\simeq\)};}}{Q}},\overset{\inserted{k^{\Lambda}}}{\overset{\scriptscriptstyle\Vert}{R}}$ are free, \deleted{the \unsure{pairings|maps} \tikz[baseline=(s.base)]{\node[inner sep=1pt](s){continuous};\node[anchor=south west,inner sep=1pt] at ([xshift=-9pt,yshift=8pt]s.north west){\inserted{bilinear}};\draw[->] ([xshift=-5pt,yshift=-2pt]s.south west) to[out=80,in=220] ([xshift=14pt,yshift=20pt]s.north west);} $\check P\times\check Q\unsure{\longrightarrow|-}\check R=\underset{\leftarrow}{\lim}_{\unsure{\alpha|\lambda}}\deleted{\unsure{R_{I}}}R_{\unsure{\Lambda_{\alpha}|\Lambda_2}}$.}
% loc 9–16, multiline Bil and Hom display
\[
\begin{aligned}
\unsure{\mathrm{bil}|\mathrm{Bil}}\,\tikz[baseline=(c.base)]{\node[inner sep=0pt](c){\(\mathrm{cont}_{k}\)};\fill ([xshift=1pt,yshift=1pt]c.south west) circle (.4pt);}\,\check P,\check Q;\check R)
 &\simeq \mathrm{Bil}\,\mathrm{cont}_{k}(\check P,\check Q;\underset{\leftarrow}{\lim}_{\gamma}(\check R_{\gamma}))\\
 &\simeq \underset{\leftarrow}{\lim}_{\gamma}\,[\mathrm{bil}\,\mathrm{cont}_{k}(\check P,\check Q;\check R_{\gamma})]\\
 &\simeq \underset{\leftarrow}{\lim}_{\gamma}\,\bigl[\underset{\mathrel{\tikz[baseline=-.3ex,x=1em,y=1em]{\draw[line width=.45pt] (0,0)--(1,0) (1,-.12)--(1,.12);}}}{\lim}_{\alpha,\beta}\,
 \mathrm{Bi}\underbrace{\mathrm{l}_{k}(\check P_{\alpha},\check Q_{\beta};\check R_{\gamma})}_{\mathrm{Hom}_{k}(\check P_{\alpha}\otimes\check Q_{\beta};\check R_{\gamma})\,]}\bigr]\\
 &\deleted{\simeq}\\
 &\simeq \underset{\leftarrow}{\lim}_{\gamma}\,\underset{\mathrel{\tikz[baseline=-.3ex,x=1em,y=1em]{\draw[line width=.45pt] (0,0)--(1,.05)--(1,-.10);}}}{\lim}_{\alpha,\beta}\,\mathrm{Hom}_{k}(R_{\gamma},P_{\alpha}\otimes Q_{\beta})\\
 &\simeq \underset{\leftarrow}{\lim}_{\gamma}\,\mathrm{Hom}_{k}(R_{\gamma},P\otimes Q)\\
 &\simeq \ldotp\,\mathrm{Hom}_{k}(R,P\otimes Q)\qquad !\quad ]
\end{aligned}
\]

% loc 17–21, We are in
\margin{\begin{tikzpicture}[x=1mm,y=1mm]\fill (0,0)--(.8,.08)--(1,-.1)--(.95,-.4)--(.35,-.5)--(.05,-.7)--cycle;\end{tikzpicture}}
\deleted{We also have the same result} \inserted{We are in the case envisaged} if $A$ is a \uline{countable} \deleted{\unsure{projective}}inductive limit of projective submodules \tikz[baseline=(s.base)]{\node[inner sep=1pt](s){, with};\node[anchor=south west,inner sep=1pt] at ([xshift=-2pt,yshift=8pt]s.north west){\inserted{of finite type}};\draw ([xshift=-3pt,yshift=-2pt]s.south west) to[out=85,in=175] ([xshift=1pt,yshift=7pt]s.north west) to[out=-5,in=200] ([xshift=14pt,yshift=14pt]s.north east);} transition morphisms \deleted{having} \deleted{complements} universal monomorphisms [i.e. having left inverses] But in the general case, we cannot forget $V(A)$ \tikz[baseline=(s.base)]{\node[inner sep=1pt](s){in favor};\node[anchor=south west,inner sep=1pt] at ([xshift=-5pt,yshift=7pt]s.north west){\inserted{and its operations}};\draw ([xshift=-3pt,yshift=-2pt]s.south west) to[out=90,in=180] ([xshift=4pt,yshift=6pt]s.north west) to[out=5,in=220] ([xshift=14pt,yshift=15pt]s.north east);} of $\check A=V(A)(k)$ $\ldotp\ldotp\ldotp\ldotp$

% loc 22–27, Quotient module and quotient map
Quotient module and submodule stable under a coaction of the coalgebra $A$ : Let $M\overset{\unsure{g|f}}{\mathrel{\tikz[baseline=-.3ex,x=1em,y=1em]{\draw[line width=.45pt] (0,0)--(1,0) (1,-.12)--(1,.12);}}}N'$ be an epimorphism, then $V(N)\subset V(M)$ monomorphism, and it is clear what it means for $V(N)$ to be stable ; we shall then say that $N$ is stable under the coaction of $N$. This also means that $\exists$ a homomorphism $N'\overset{\omega_{N'}}{\mathrel{\tikz[baseline=-.3ex,x=1em,y=1em]{\draw[line width=.45pt] (0,0)--(1,0) (1,-.12)--(1,.12);}}}N'\otimes A$ making commutative
% loc 27–28, commutative square
\diagram{\begin{tikzcd}[ampersand replacement=\&,column sep=large,execute at end picture={\draw ([xshift=-3pt,yshift=-2pt]nout.west)--([xshift=-3pt,yshift=2pt]nout.west);}]
M\arrow[r,"\omega_M"]\arrow[d] \& M\otimes A\arrow[d,"g\otimes\mathrm{id}_{A}"]\\
N'\arrow[r,dash,"\omega_{N'}"] \& |[alias=nout]| N'\otimes A
\end{tikzcd}}
% loc 29–30, last two lines
and this $\omega_{N}$ is necessarily unique and satisfies the unit and associativity axioms. If $A$ is \unsure{an affine group|an affine group|an affine algebra}
% archive sheet number: 27
